% Chapter 3: Partial Differential Equations as Solution Operators
% Part I of the book.

\chapter{Partial Differential Equations as Solution Operators}
\label{ch:part1-partial-differential-equations-as-solution}
\label{sec:pde_operators}
% -----------------------------------------------------------------------------

A partial differential equation (PDE) does not, by itself, pose a learning
problem. What a surrogate is asked to reproduce is not the residual but the
consequence
of the residual's vanishing: the field selected by the prescribed data. Write a
model problem as $\mathcal{P}(u; a) = 0$, where $\mathcal{P} : \mathcal{U} \times
\mathcal{A} \to \mathcal{Z}$ is a generally nonlinear differential expression
taking values in a residual space $\mathcal{Z}$, $a \in \mathcal{A}$ denotes the
input parameters --- an initial state, a boundary trace, a forcing term, or a
spatially varying coefficient --- and $u \in \mathcal{U}$ denotes the solution.
The object one actually samples and approximates is then not $\mathcal{P}$ but the
\index{solution operator}\textbf{solution
operator} $\mathcal{G} : \mathcal{A} \to \mathcal{U}$ carrying each admissible
datum to the solution it determines. The differential expression is a
verification device; the operator is the target. Spaces, norms, and duality
pairings follow the conventions fixed in
\chapref{ch:part1-function-spaces-duality-and-sobolev-regula} and collected in
Table~\ref{tab:notation-master}.

Recasting a residual as a map is where the mathematics
enters. Three questions must be answered before $\mathcal{G}$ exists
as an object of study: does every admissible input admit at least one solution,
is that solution unique, and does it depend continuously on the input --- small
perturbations of $a$ in the norm of $\mathcal{A}$ producing small changes of
$\mathcal{G}(a)$ in the norm of $\mathcal{U}$? These are the three
\index{well-posedness}\index{Hadamard well-posedness}conditions of Hadamard
well-posedness, and each answer is a statement about a specific \emph{pair} of
norms, never about the equation in the abstract. The same evolution problem is
Lipschitz between one choice of spaces and discontinuous between another, and
the choice is not free: it is dictated by whatever the data of the physical
problem naturally control. A second distinction, which must be kept precise,
runs through every example below. A \index{solution operator!forward (time-evolution)}forward
solution operator maps a state at one time to a state at a later time, so
domain and codomain coincide; a
\index{coefficient-to-state map}\textbf{coefficient-to-state map} takes a
parameter field as input and a state as output, so the two spaces differ and the
map is typically smoothing. The two impose different demands on an architecture.

We treat three families in turn: the elliptic Darcy problem, which supplies the
cleanest complete answer; the Green-kernel representation, which exposes why
solution operators are nonlocal; and the parabolic and hyperbolic evolution
problems, where the answer deteriorates as viscosity vanishes. Every estimate
derived here becomes a hypothesis in \chapref{ch:part1-operator-approximation-theory-and-universa},
so constants and spaces are recorded with intent.

\section{Elliptic Problems and Darcy Flow}
\index{Darcy flow}Stationary flow of a single phase through a saturated porous
medium is described by two relations: mass conservation, $\nabla \cdot \mathbf{q} = f$,
and Darcy's empirical law for the flux, $\mathbf{q} = -a \nabla u$, where $u$ is
the pressure potential and $a$ the permeability. Eliminating $\mathbf{q}$ yields
the steady-state diffusion problem on a bounded domain
$\Omega \subset \mathbb{R}^d$,
\begin{equation}
  \label{eq:darcy}
  \begin{cases}
    -\nabla \cdot \big( a(x) \nabla u(x) \big) = f(x), & x \in \Omega, \\
    u(x) = 0, & x \in \partial\Omega.
  \end{cases}
\end{equation}
Here $a \in \mathcal{A} = L^\infty(\Omega)$ is the permeability coefficient and
$f \in H^{-1}(\Omega)$ is a fixed source. Two different maps live on
\eqref{eq:darcy}, and conflating them is a common error. Holding $a$ fixed, the
forcing-to-state map $f \mapsto u$ is \emph{linear}, because the differential
expression is linear in $u$; superposition holds and the representation is the
Green kernel of the next section. Holding $f$ fixed and letting $a$ vary --- the
problem this book takes as its elliptic model and the benchmark used throughout
\partref{part:2} --- the coefficient-to-state map $a \mapsto u$ is genuinely
\emph{nonlinear}, because $a$ multiplies $\nabla u$ inside the divergence
operator. Nothing in \eqref{eq:darcy} linearizes the second map.

A classical solution of \eqref{eq:darcy} is unavailable when $a$ is merely
essentially bounded: $\nabla \cdot (a \nabla u)$ need not exist pointwise. The
standard remedy, and the one that makes the estimate below transparent, is to
test against $v \in H^1_0(\Omega)$ and integrate by parts, the boundary term
dropping out because the test functions vanish on $\partial\Omega$. This yields
the weak formulation: find $u \in H^1_0(\Omega)$ such that
\begin{equation}
  \label{eq:ch03-weakform}
  B_a(u, v) \coloneqq \int_\Omega a(x) \nabla u \cdot \nabla v \, dx = \langle f, v \rangle_{H^{-1}, H^1_0} \quad \forall v \in H^1_0(\Omega).
\end{equation}
The bilinear form in \eqref{eq:ch03-weakform} is well defined on
$H^1_0 \times H^1_0$ since $a \in L^\infty$ and $\nabla u, \nabla v \in L^2$;
it is bounded for every such $a$, but coercive only once the coefficient is
bounded away from zero, which is the content of the next hypothesis.

\begin{assumption}[Uniform Ellipticity]
\label{ass:ellipticity}
There exist constants $0 < a_{\min} \le a_{\max} < \infty$ such that the
admissible coefficient class is
\begin{equation}
  \mathcal{A}_{\text{adm}} = \left\{ a \in L^\infty(\Omega) : a_{\min} \le a(x) \le a_{\max} \quad \text{a.e. } x \in \Omega \right\}.
  \label{eq:ch03-adm}
\end{equation}
\end{assumption}

\index{uniform ellipticity}Assumption~\ref{ass:ellipticity} makes $B_a$ coercive.
For $a \in \mathcal{A}_{\text{adm}}$ given by \eqref{eq:ch03-adm}, the form
satisfies the pair of estimates
\begin{equation}
  B_a(v,v) \ge a_{\min} \|\nabla v\|_{L^2}^2
  \quad \text{and} \quad
  |B_a(w, v)| \le a_{\max} \|\nabla w\|_{L^2} \|\nabla v\|_{L^2},
  \label{eq:ch03-coercive}
\end{equation}
that is, coercivity with constant $a_{\min}$ and continuity with constant
$a_{\max}$. Both follow from moving the pointwise bounds on $a$ outside the
integral in \eqref{eq:ch03-weakform}, and the coercive half of
\eqref{eq:ch03-coercive} is precisely what fails if $a$ is allowed to approach
zero on a set of positive measure. On $H^1_0(\Omega)$ the Poincar{\'e}
\index{Poincar{\'e} inequality}inequality makes
$\|\nabla \cdot\|_{L^2}$ an equivalent norm, and the right-hand side of
\eqref{eq:ch03-weakform} is continuous by definition of $H^{-1}(\Omega)$, the
dual space introduced in \chapref{ch:part1-function-spaces-duality-and-sobolev-regula}.
The Lax--Milgram \index{Lax--Milgram theorem}\index{weak formulation}theorem~\cite{brezis2011functional, evans2010partial}
then supplies existence, uniqueness, and the a priori bound recorded in the next
proposition, all in one step.

\begin{proposition}[Properties of the Darcy Solution Operator]
\index{Darcy flow!solution operator}Under Assumption~\ref{ass:ellipticity}, for
every $a \in \mathcal{A}_{\text{adm}}$ and $f \in H^{-1}(\Omega)$, the
Lax--Milgram theorem guarantees a unique weak solution $u \in H^1_0(\Omega)$
satisfying \eqref{eq:ch03-weakform}. The solution operator
$\mathcal{G} : \mathcal{A}_{\text{adm}} \subset L^\infty(\Omega) \to H^1_0(\Omega)$,
$a \mapsto u$, has three properties of distinct character. Its first property is
a \textbf{size} statement, \index{uniform boundedness}global boundedness:
\begin{equation}
  \| \mathcal{G}(a) \|_{H^1_0(\Omega)} \le \frac{1}{a_{\min}} \|f\|_{H^{-1}(\Omega)}
  \quad \text{for all } a \in \mathcal{A}_{\text{adm}}.
  \label{eq:ch03-bounded}
\end{equation}
Its second property is a \textbf{stability} statement, \index{Lipschitz continuity!of a solution operator}Lipschitz continuity: for any
$a_1, a_2 \in \mathcal{A}_{\text{adm}}$,
\begin{equation}
  \| \mathcal{G}(a_1) - \mathcal{G}(a_2) \|_{H^1_0(\Omega)} \le \frac{\|f\|_{H^{-1}(\Omega)}}{a_{\min}^2} \|a_1 - a_2\|_{L^\infty(\Omega)}.
  \label{eq:ch03-lipschitz}
\end{equation}
Its third property is a \textbf{smoothness} statement, \index{Fr{\'e}chet differentiability}Fr{\'e}chet differentiability: $\mathcal{G}$ is infinitely
Fr{\'e}chet differentiable on the interior of $\mathcal{A}_{\text{adm}}$, with
the interior taken relative to the $L^\infty$ norm~\cite{stuart2010inverse}. The
three should not be read as redundancies. The size estimate
\eqref{eq:ch03-bounded} says the image of the admissible set is a bounded subset
of $H^1_0$; Lipschitz continuity says the
parameter-to-solution correspondence is stable, which is what a regression loss
can hope to reproduce; differentiability says the map has a well-defined linear
response to coefficient perturbations, the fact that inverse problems and
gradient-based design rest on.
\end{proposition}

Two consequences of \eqref{eq:ch03-lipschitz} matter for what follows. First,
the Lipschitz constant degrades quadratically as ellipticity degenerates:
$a_{\min}^{-2}$ diverges for high-contrast media with near-vacuum layers, so the
problem becomes stiff exactly where the coefficient is most interesting. Second,
the input metric is $L^\infty$, the strongest natural one, and it overstates the
sensitivity that matters for training. Perturbing $a$ on a set $E$ of small
measure leaves $\|a_1 - a_2\|_{L^\infty}$ of order one, so
\eqref{eq:ch03-lipschitz} charges the full amount, yet the argument that
produced the bound localizes the effect: the difference identity
\eqref{eq:ch03-subtract}, derived in the proof sketch below, has a right-hand
side that couples to $\nabla u_2$ only on $E$, and the resulting
estimate $\|\nabla(u_1-u_2)\|_{L^2} \le a_{\min}^{-1} \|a_1-a_2\|_{L^\infty}
\|\nabla u_2\|_{L^2(E)}$ tends to zero with the measure of $E$. Running the
estimate with $a_1 - a_2 \in L^p$ and $\nabla u_2 \in L^q$, where
$1/p + 1/q = 1/2$, makes the same point in norm form, at the cost of an
integrability assumption on the gradient that Assumption~\ref{ass:ellipticity}
alone does not deliver.

\begin{corollary}[Precompactness of the Darcy solution set]
\label{cor:ch03-precompact}
Under Assumption~\ref{ass:ellipticity}, the image
$\mathcal{G}(\mathcal{A}_{\text{adm}})$ is bounded in $H^1_0(\Omega)$ and
therefore precompact in $L^2(\Omega)$.
\end{corollary}

Corollary~\ref{cor:ch03-precompact} is immediate: the bound
\eqref{eq:ch03-bounded} gives boundedness in $H^1_0$, and the
Rellich--Kondrachov compact embedding
\index{Rellich--Kondrachov compact embedding!consequence for solution sets}of
\chapref{ch:part1-function-spaces-duality-and-sobolev-regula} does the rest. It
is also the first place the topological obstruction of
\chapref{ch:part1-function-spaces-duality-and-sobolev-regula} is discharged: in
infinite dimensions bounded does not imply compact, but bounded in a
\emph{stronger} space does. Uniform approximation of $\mathcal{G}$ will be
possible precisely because the sets over which it is claimed are small in this
sense.

\begin{proof}[Proof Sketch for Lipschitz Continuity]
Let $u_1 = \mathcal{G}(a_1)$ and $u_2 = \mathcal{G}(a_2)$. Subtracting the weak
forms \eqref{eq:ch03-weakform} yields
\begin{equation}
  \int_\Omega a_1 \nabla (u_1 - u_2) \cdot \nabla v \, dx = \int_\Omega (a_2 - a_1) \nabla u_2 \cdot \nabla v \, dx \quad \forall v \in H^1_0(\Omega).
  \label{eq:ch03-subtract}
\end{equation}
Setting $v = u_1 - u_2$ in \eqref{eq:ch03-subtract} --- legal because the
difference of two $H^1_0$
functions is itself an admissible test function --- and invoking uniform
ellipticity on the left-hand side and H{\"o}lder's inequality on the right gives
\begin{equation}
  a_{\min} \|\nabla(u_1 - u_2)\|_{L^2}^2 \le \|a_1 - a_2\|_{L^\infty} \|\nabla u_2\|_{L^2} \|\nabla(u_1 - u_2)\|_{L^2}.
  \label{eq:ch03-energy}
\end{equation}
Dividing \eqref{eq:ch03-energy} by $\|\nabla(u_1 - u_2)\|_{L^2}$ and inserting
the a priori bound $\|\nabla u_2\|_{L^2} \le a_{\min}^{-1} \|f\|_{H^{-1}}$,
which is \eqref{eq:ch03-bounded} in gradient form, confirms
\eqref{eq:ch03-lipschitz}.
\end{proof}

\begin{remark}[Which norm, and on which ensembles]
The bound \eqref{eq:ch03-bounded} is written
with $\|w\|_{H^1_0} = \|\nabla w\|_{L^2}$; with the full norm
$(\|w\|_{L^2}^2 + \|\nabla w\|_{L^2}^2)^{1/2}$ an additional Poincar{\'e}
constant appears on a bounded domain. The distinction changes constants, not
conclusions. More consequential is the mismatch between
Assumption~\ref{ass:ellipticity} and typical training ensembles. The Darcy
benchmark of \partref{part:2} draws permeabilities from a log-normal random
field, whose samples are strictly positive almost surely but satisfy
\eqref{eq:ch03-adm} with \emph{sample-dependent} values of $a_{\min}$. Then
\eqref{eq:ch03-lipschitz} holds path by path while the Lipschitz constant is a
heavy-tailed random variable, and no uniform bound exists over the ensemble
unless the sampled field is truncated. Truncation is the practical remedy; it
also restores the two-sided bounds in \eqref{eq:ch03-adm} in a deterministic
sense, at the price of altering the input measure.
\end{remark}

\section{Green's Functions and Nonlocal Integral Operators}
Fix the coefficient. When $\mathcal{L}_a = -\nabla \cdot (a \nabla)$ is linear
and invertible, its inverse is represented by an integral operator:
\begin{equation}
  u(x) = (\mathcal{L}_a^{-1} f)(x) = \int_\Omega G_a(x, y) f(y) \, dy,
  \label{eq:ch03-green}
\end{equation}
where
\index{Green's function}$G_a : \Omega \times \Omega \to \mathbb{R} \cup \{\infty\}$
is the \textbf{Green's function} of the boundary value problem.
Equation~\eqref{eq:ch03-green} makes three things visible at once. It exhibits
the solution as a \emph{nonlocal} average: the value at $x$ responds to forcing
applied anywhere in $\Omega$, with weight $G_a(x,y)$. It exposes the singularity
structure, since $G_a$ diverges on the diagonal $x = y$ --- like
$|x-y|^{2-d}$ for $d \ge 3$ and like $\log|x-y|$ for $d=2$, the fundamental
solution of the Laplacian being the leading term --- yet remains integrable,
which is exactly what the weak formulation encodes. And it records the
\emph{regularity gain}: for smooth coefficients on smooth or convex domains,
classical elliptic regularity turns \eqref{eq:ch03-green} into
$\|u\|_{H^2(\Omega)} \le C\|f\|_{L^2(\Omega)}$, so $\mathcal{L}_a^{-1}$ recovers
two derivatives~\cite{evans2010partial, brezis2011functional}. The gain is what
Corollary~\ref{cor:ch03-precompact} exploited, and it is the elliptic analogue of
the parabolic smoothing discussed below.

For a fixed coefficient there is a second representation, closer to the geometry
of the operator than the kernel itself. The Dirichlet realization of
$\mathcal{L}_a$ is self-adjoint --- the coefficient is scalar and the operator in
divergence form --- and has compact resolvent on $L^2(\Omega)$, again by
Rellich--Kondrachov. It therefore possesses an orthonormal eigenbasis
$\{\phi_j\}_{j \ge 1}$ with eigenvalues $0 < \lambda_1 \le \lambda_2 \le \dots
\uparrow \infty$, and \eqref{eq:ch03-green} expands as
\begin{equation}
  G_a(x, y) = \sum_{j=1}^{\infty} \frac{\phi_j(x)\, \phi_j(y)}{\lambda_j},
  \qquad
  u(x) = \sum_{j=1}^{\infty} \frac{\langle f, \phi_j \rangle_{L^2}}{\lambda_j}\, \phi_j(x).
  \label{eq:ch03-eigen}
\end{equation}
The series in \eqref{eq:ch03-eigen} is the precise sense in which solving
elliptically means dividing by the spectrum: each mode of the data is amplified
by $1/\lambda_j$, and since the Hilbert--Schmidt norm of the kernel is
$(\sum_j \lambda_j^{-2})^{1/2}$, Weyl's law $\lambda_j \sim c\, j^{2/d}$ makes
the amplification summable fast enough that $G_a \in L^2(\Omega \times \Omega)$
for
$d \le 3$. Two lessons follow. The first is that keeping the leading eigenpairs
is the best low-rank linear approximation of $\mathcal{L}_a^{-1}$ in operator
norm, which is why basis truncation is the first move of any spectral method. The
second is that the eigenpairs themselves depend
on $a$, so the basis in \eqref{eq:ch03-eigen} is not fixed in advance but moves
with every sampled coefficient field --- the difficulty a learned kernel is meant
to absorb.

The mechanism is completely explicit in one case that recurs throughout this
book. For the constant-coefficient Laplacian on the periodic torus
$\mathbb{T}^d$, translation invariance forces $G_a(x,y) = G(x-y)$, and the
Fourier basis diagonalizes the kernel: for $-\Delta u = f$ with
$\int_{\mathbb{T}^d} f\,dx = 0$ and $u$ of mean zero,
\begin{equation}
  \hat{u}(k) = \frac{\hat{f}(k)}{4\pi^2 |k|^2}, \qquad k \in \mathbb{Z}^d \setminus \{0\},
  \label{eq:ch03-poissonfourier}
\end{equation}
which is \eqref{eq:ch03-green} read in the Fourier basis, the one representation
in which translation invariance turns the kernel into a function of the
difference $x - y$ and the solution operator into a multiplier. The multiplier in
\eqref{eq:ch03-poissonfourier} decays like $|k|^{-2}$, so low frequencies are
amplified and high frequencies are suppressed; the operator is compact on $L^2$
and its output spectrum decays two powers faster than its input. That single
observation is the seed of \chapref{ch:part1-the-fourier-neural-operator-spectral-theor},
where a learned kernel is deliberately restricted to exactly this
translation-invariant, Fourier-diagonal form so that the dense
$\mathcal{O}(N^2)$ evaluation of \eqref{eq:ch03-green} collapses to
$\mathcal{O}(N\log N)$. The trade is honest: stationarity discards the
$x$-and-$y$ dependence that a heterogeneous coefficient field in
\eqref{eq:darcy} genuinely requires.

Nonlocality is the property an architecture must reproduce, and it is the reason
local stencil methods need depth to represent elliptic physics. A convolution
with a fixed finite kernel cannot transmit information from the far side of the
domain to a point in one layer, whereas \eqref{eq:ch03-green} does so by
construction. This is the thread \chapref{ch:part1-discretization-invariance-and-continuum-li}
picks up when it asks which discrete stencils converge to a continuum integral
at all.

Time dependence promotes the kernel to a causal one. On $\mathbb{T}$, the heat
semigroup $S(t) = e^{\nu t \partial_{xx}}$ has Fourier symbol
$e^{-\nu (2\pi k)^2 t}$, and the semilinear viscous Burgers equation admits the
mild, or variation-of-constants, representation
\begin{equation}
  u(\cdot, t) = S(t) u_0 - \int_0^t S(t-s)\, \partial_x\!\left(\tfrac{1}{2} u(\cdot, s)^2\right) ds,
  \qquad S(t) = e^{\nu t \partial_{xx}}.
  \label{eq:ch03-duhamel}
\end{equation}
Equation~\eqref{eq:ch03-duhamel} is the closest nonlinear analogue of a Green
representation, and the difference is instructive: the integral is against a
known linear kernel, but the integrand contains the unknown solution itself.
Superposition fails for nonlinear problems, so no kernel in $(x,y)$ alone
represents $\mathcal{G}$; only fixed-point arguments recover existence. A neural
operator layer built from integral kernels therefore borrows the \emph{shape} of
\eqref{eq:ch03-green}, not its \emph{guarantee}.

\takeaway{The Architectural Inspiration for Neural Operators}{The solution to any linear elliptic or parabolic PDE can be expressed as a continuous integral transform with kernel $G_a(x,y)$. The generic neural operator layer directly parameterizes a trainable, nonlinear surrogate of this integral kernel representation.}

\section{Parabolic and Nonlinear Hyperbolic Evolution}
Evolution problems are the forward operators announced at the head of this chapter.
For the nonlinear viscous Burgers' equation on $\mathbb{T}$,
\index{Burgers equation}\index{Burgers equation!viscous}
\begin{equation}
  \partial_t u + u \partial_x u = \nu \partial_{xx} u, \qquad u(x, 0) = u_0(x),
  \label{eq:ch03-burgers}
\end{equation}
the solution operator $\mathcal{G}_t : u_0 \mapsto u(\cdot, t)$ maps initial
conditions $u_0 \in L^2(\mathbb{T})$ to smooth states $u(\cdot, t) \in H^s(\mathbb{T})$
for any $s > 0$ when $\nu > 0$. Its argument is a state, not a coefficient, so
domain and codomain are function spaces of the same kind. Because
\eqref{eq:ch03-burgers} is autonomous,
the family $\{\mathcal{G}_t\}_{t \ge 0}$ is a semigroup,
\index{semigroup property}
\begin{equation}
  \mathcal{G}_{t+\tau} = \mathcal{G}_t \circ \mathcal{G}_\tau, \qquad \mathcal{G}_0 = \mathrm{Id},
  \label{eq:ch03-semigroup}
\end{equation}
which is the structural reason a single learned time-step map can be composed
recursively --- and the reason such rollouts inherit and amplify the error of
one step. The smoothing is best read through the slope $w = \partial_x u$:
differentiating \eqref{eq:ch03-burgers} gives
\begin{equation}
  \partial_t w + u\, \partial_x w = \nu\, \partial_{xx} w - w^2,
  \qquad w = \partial_x u,
  \label{eq:ch03-gradient}
\end{equation}
whose reaction term is signed. Since $w$ is periodic it attains a maximum, and at
a positive maximum one has $\partial_x w = 0$ and $\partial_{xx} w \le 0$, so
\eqref{eq:ch03-gradient} leaves $d(\sup_x w)/dt \le -(\sup_x w)^2$: upward slopes
cannot grow, and $\sup_x \partial_x u(\cdot, t) \le \sup_x u_0'(x)$ for as long
as the solution stays smooth. At a negative minimum the two terms on the right
of \eqref{eq:ch03-gradient} carry opposite signs and compete, the reaction
$-w^2$ deepening the downswing while diffusion flattens it. That contest is the
shock, and \index{viscosity}viscosity is what keeps $\inf_x w$ finite. Without
viscosity the protection goes: along characteristics of the inviscid equation $w$
obeys $dw/dt = -w^2$,
whence $w(t) = w_0/(1 + t w_0)$ and the first gradient catastrophe occurs at
$t_* = \big(\max_x[-u_0'(x)]\big)^{-1}$ wherever the initial slope is negative.
For data of the form $u_0(x) = A\sin(2\pi x)$ this gives $t_* = (2\pi A)^{-1}$,
so amplitudes of order unity place the catastrophe well inside the horizon
$t \in [0,1]$ of the Burgers benchmark in \partref{part:2}, whose steepening is
this mechanism at small positive $\nu$.

At positive $\nu$ the singularity becomes a transition layer. A compressive
shock connecting states $u^- > u^+$ travels at the Rankine--Hugoniot speed
$s = (u^+ + u^-)/2$ and has the viscous profile
\index{shock wave}\index{Rankine--Hugoniot condition}
\begin{equation}
  \bar{u}(\xi) = \frac{u^- + u^+}{2} - \frac{u^- - u^+}{2}
  \tanh\!\left(\frac{(u^- - u^+)\,\xi}{4\nu}\right),
  \qquad \xi = x - s t,
  \label{eq:ch03-shockprofile}
\end{equation}
so the layer thickness scales like $4\nu/(u^- - u^+)$
\cite{leveque2002finite}. As $\nu \to 0^+$, \eqref{eq:ch03-shockprofile} tends to
a jump, shocks steepen, and energy is transferred to
arbitrarily high spatial frequencies. Two distinct losses accompany that limit.
The first is regularity. At a time when the inviscid limit has developed a
jump, the size of $\mathcal{G}_t^\nu$ --- the Burgers operator at viscosity
$\nu$ --- as a map
$L^2(\mathbb{T}) \to
H^s(\mathbb{T})$ diverges as $\nu \to 0^+$ for every $s > \tfrac{1}{2}$;
since the map is nonlinear, size is measured on a ball $\|u_0\|_{L^2} \le R$
by $\sup \|\mathcal{G}_t^\nu(u_0)\|_{H^s(\mathbb{T})}$, and the supremum diverges
once the ball contains data whose inviscid limit jumps. A bound uniform in $\nu$
would give a subsequence converging weakly in $H^s$ and strongly in $L^2$ to that
inviscid limit, an $H^s$ function with a jump --- yet a function with a jump
belongs to
$H^s(\mathbb{T})$ only for $s < \tfrac{1}{2}$. The second is the quality of the
stability estimate: the metric in which the inviscid theory is well posed is
$L^1$, where the entropy semigroup is nonexpansive \cite{leveque2002finite}, and
an $L^1$ contraction controls no Sobolev seminorm. Selecting the entropy solution
\index{entropy solution}--- the vanishing viscosity limit singles it out --- is
therefore part of specifying the operator, not an
implementation detail. This is why \eqref{eq:ch03-burgers} is an essential
benchmark for any architecture meant to capture sharp multiscale gradients, and
why the smoothing above carries no uniformity
as $\nu \to 0^+$: the factor $e^{-\nu (2\pi k)^2 t}$ in the mild
form \eqref{eq:ch03-duhamel} is where dissipation and cascade trade off.

For the two-dimensional incompressible Navier--Stokes equations in
vorticity--streamfunction formulation, \index{Navier--Stokes equations!vorticity-streamfunction formulation}
\begin{equation}
  \partial_t \omega + \mathbf{u} \cdot \nabla \omega = \nu \Delta \omega + f, \qquad \nabla \cdot \mathbf{u} = 0, \quad \omega = \nabla \times \mathbf{u},
  \label{eq:ch03-navierstokes}
\end{equation}
the solution operator maps $\omega(\cdot, 0) \in L^2(\mathbb{T}^2)$ to
$\omega(\cdot, T) \in L^2(\mathbb{T}^2)$. Two structural facts distinguish
\eqref{eq:ch03-navierstokes} from \eqref{eq:ch03-burgers}. First, the velocity is
recovered from vorticity by the Biot--Savart law,
$\mathbf{u} = \nabla^{\perp}(-\Delta_{\mathbb{T}^2})^{-1}\omega$, a singular
integral of order zero on the torus, so the nonlinearity is of pure transport
type and $\|\nabla \mathbf{u}\|_{L^2} = \|\omega\|_{L^2}$. Second, incompressibility
kills the nonlinear term in the energy balance: testing \eqref{eq:ch03-navierstokes}
against $\omega$ gives
\index{Biot--Savart law}\index{enstrophy}
\begin{equation}
  \frac{1}{2}\frac{d}{dt}\|\omega(\cdot, t)\|_{L^2}^2
  = -\nu \|\nabla \omega(\cdot, t)\|_{L^2}^2 + \langle f, \omega \rangle_{L^2},
  \label{eq:ch03-enstrophy}
\end{equation}
because $\int_{\mathbb{T}^2} (\mathbf{u}\cdot\nabla\omega)\,\omega\,dx =
\tfrac{1}{2}\int_{\mathbb{T}^2} \mathbf{u}\cdot\nabla(\omega^2)\,dx = 0$. Once the
forcing is bounded by $\|f\|_{L^2}\|\omega\|_{L^2}$ and the inequality is closed
with Gr\"onwall's lemma, the balance \eqref{eq:ch03-enstrophy} is a closed
\emph{a priori} estimate in two space dimensions --- the reason the forward
operator in \eqref{eq:ch03-navierstokes} is
globally defined; in three dimensions vortex stretching adds a term of this same
sign and the estimate no longer closes, so nothing here is claimed beyond two
dimensions. Long-time behavior
supplies one further structural fact: the forced two-dimensional dynamics lie on
a finite-dimensional global attractor \index{global attractor}\index{dissipative system}, implying that
the nonlinear solution operator maps bounded sets in $L^2$ into compact, smooth
subsets. We take that statement as a hypothesis, not a result proved here:
it is exactly the kind of compactness that makes uniform approximation
conceivable, and its consequences for long autoregressive rollouts are confronted
in \chapref{ch:part1-open-theoretical-problems-and-mathematical}.

Taken together, the examples fix the vocabulary the rest of
\partref{part:1} uses. Elliptic
coefficient-to-state maps such as $\mathcal{G}$ in \eqref{eq:ch03-weakform} are
bounded, Lipschitz, smooth, and regularizing between different spaces; forward
operators such as $\mathcal{G}_t$ in \eqref{eq:ch03-semigroup} are state-to-state
semigroups, smoothing under viscosity and contractive only in a weaker norm in
the hyperbolic limit. Both are continuous maps between infinite-dimensional
spaces, which is the class \index{operator learning}neural operators approximate,
and neither is a learning target until a measure on $\mathcal{A}$ is chosen; the
admissible set \eqref{eq:ch03-adm}, moreover, is far from all of $L^\infty$.

Well-posedness gives us a map, estimates on it, and a topology in which its
image is small. None of that says the map is \emph{representable}. We have not
exhibited a single finite-dimensional recipe for $\mathcal{G}$, have not said how
much accuracy a given number of parameters buys, and have not distinguished the
difficulty of approximating the linear Green operator of \eqref{eq:ch03-green}
from that of approximating the Burgers semigroup of \eqref{eq:ch03-burgers}. Those
are questions of approximation theory, and the answer depends on precisely the
properties catalogued above: continuity of the target, boundedness and
precompactness of the sets over which error is measured --- the role
Corollary~\ref{cor:ch03-precompact} plays --- and the regularity gain that fixes
how quickly an optimally truncated basis loses accuracy. \chapref{ch:part1-operator-approximation-theory-and-universa}
takes these as data. It asks whether networks of finite width and depth can
approximate continuous operators between function spaces to arbitrary precision,
in what sense the classical result of Chen and Chen is an existence statement
about sensor evaluations and output bases, and how the branch--trunk and
integral-kernel architectures that followed turn that statement into a trainable
model.
